\exercisesection{Banach 空间自同态的谱}

在下面的习题中，所有 Banach 空间都取在 C 上。

\begin{enumerate}
\def\labelenumi{\arabic{enumi})}
\tightlist
\item
  设 \(\mathrm{E} = \mathcal{C}([0,1])\)，设 \(u\) 为由下式定义的 \(\mathrm{E}\) 的自同态
\end{enumerate}

\[
u (f) (t) = \int_ {0} ^ {t} f (t) d t.
\]

\begin{enumerate}
\def\labelenumi{\alph{enumi})}
\item
  \(u\) 的谱半径为零。（估计 \(u^n\) 的范数。）
\item
  \(u\) 的谱仅由 0 组成，且对每个 \(n \geqslant 1\)，\(u^{n}\) 的核仅由 0 组成。
\end{enumerate}

（特别地，0 是 \(u\) 的谱的一个孤立点，但诸 \(u^n\)（\(n \geqslant 1\)）的核的并集的闭包并不等于谱子空间 \(\mathrm{E}_0(x)\)。）

\begin{enumerate}
\def\labelenumi{\arabic{enumi})}
\setcounter{enumi}{1}
\item
  存在一个 Banach 空间 E 及 E 的一个自同态 u，使得 0 是 u 的特征值，但 0 在 \(\mathrm{Sp}(u)\) 中不是孤立点。
\item
  给出一个 Hilbert 空间 E 与一个自同态 \(u \in \mathcal{L}(\mathrm{E})\) 的例子，使得 0 是 u 的特征值，但不是 \(u^{*}\) 的特征值。
\item
  设 \(E = \ell^{2}(\mathbf{Z})\)，设 \(u\) 为满足 \(u((\lambda_{n})) = (\lambda_{n+1})\) 的自同态。
\end{enumerate}

\begin{enumerate}
\def\labelenumi{\alph{enumi})}
\item
  证明 \(u\) 是酉自同态。
\item
  计算 \(u\) 的谱。
\item
  \(u\) 的有限谱是空集。
\end{enumerate}

\begin{enumerate}
\def\labelenumi{\arabic{enumi})}
\setcounter{enumi}{4}
\tightlist
\item
  设 \(s\) 为 \(\ell^2 (\mathbf{N})\) 的由下式确定的自同态
\end{enumerate}

\[
s (a _ {0}, a _ {1}, \dots , a _ {n}, \dots) = (a _ {1}, \dots , a _ {n}, \dots).
\]

证明 \(s\) 的数值域是 C 中的开单位圆盘 \(\Delta\)，而 \(s\) 的谱是 C 中的闭单位圆盘。求出 \(s\) 的特征值。对 \(s^{*}\)（\(s\) 的伴随自同态）作同样的讨论。

\begin{enumerate}
\def\labelenumi{\arabic{enumi})}
\setcounter{enumi}{5}
\item
  设 \(u \in \mathcal{L}(\mathbf{C}^n)\) 为迹为零的自同态。存在 \(\mathbf{C}^n\) 的一个标准正交基 B，使得 \(u\) 在基 B 下的矩阵的所有对角元均为零。（证明 0 属于 \(u\) 的数值域。）
\item
  设 \(u\) 为复 Hilbert 空间 E 的自同态，使得 \(u\) 的谱包含于 C 中的开单位圆盘。设 F 为由 E 的元素组成的平方可和序列所构成的 Hilbert 空间 \(\ell_{\mathrm{E}}^{2}(\mathbf{N})\)（EVT，V，第 18 页，例），并设 \(s_{\mathrm{F}}\) 为 F 的由下式确定的自同态
\end{enumerate}

\[
s _ {\mathrm{F}} (a _ {0}, a _ {1}, \ldots , a _ {n}, \ldots) = (a _ {1}, \ldots , a _ {n} \ldots).
\]

存在 F 的一个闭子空间 H 与一个连续的双射线性映射 \(v: E \rightarrow H\)，使得 \(u = v^{-1}s_{F}v\)。

\begin{enumerate}
\def\labelenumi{\arabic{enumi})}
\setcounter{enumi}{7}
\tightlist
\item
  设 \(u\) 为复 Hilbert 空间 E 的一个自同态。设 \(S_u\) 为 E 的形如 \(wuw^{-1}\) 的自同态所成的集合，其中 \(w\) 取遍 E 的可逆自同态。
\end{enumerate}

\begin{enumerate}
\def\labelenumi{\alph{enumi})}
\tightlist
\item
  我们有
\end{enumerate}

\[
\varrho (u) = \inf _ {v \in \mathrm{S} _ {u}} \| v \|.
\]

（可设 \(\varrho(u) \leqslant 1\)；对自同态 \(u_{\varepsilon} = (1 + \varepsilon)^{-1} u\)（\(\varepsilon > 0\)）应用前一个习题。）

\begin{enumerate}
\def\labelenumi{\alph{enumi})}
\setcounter{enumi}{1}
\tightlist
\item
  \(u\) 的谱的凸包等于诸子集 \(\iota(v)\)（\(v \in S_u\)）的交（"Hildebrandt 定理"）。
\end{enumerate}

¶ 9) 设 E 为复 Banach 空间。

\begin{enumerate}
\def\labelenumi{\alph{enumi})}
\item
  设 \(u \in \mathcal{L}(\mathrm{E})\)。考察 \(\gamma_{u}\) 与 \(\delta_{u}\)，即 \(\mathcal{L}(\mathrm{E})\) 的由 \(\gamma_{u}(v) = u \circ v\) 与 \(\delta_{u}(v) = v \circ u\) 定义的自同态。我们有 \(\operatorname{Sp}(u) = \operatorname{Sp}(\gamma_{u}) = \operatorname{Sp}(\delta_{u})\)，其中后两个谱是相对于代数 \(\mathcal{L}(\mathcal{L}(\mathrm{E}))\) 取的。
\item
  对每个函数 \(f\)，只要它在 \(u\) 的谱的某邻域内全纯，就有 \(\boldsymbol{\gamma}_{f(u)} = f(\boldsymbol{\gamma}_{u})\) 与 \(\boldsymbol{\delta}_{f(u)} = f(\boldsymbol{\delta}_{u})\)。
\item
  设 \(u_{1}\) 与 \(u_{2}\) 属于 \(\mathcal{L}(\mathrm{E})\)。设 \(\varrho\) 为映射 \(v \mapsto u_{1} \circ v \circ u_{2}\)，它是 \(\mathcal{L}(\mathrm{E})\) 的一个自同态。它的谱等于 \(\operatorname{Sp}(u_{1})\operatorname{Sp}(u_{2}) \subset \mathbf{C}\)。
\end{enumerate}

\(d)\) 设 \(u_{1}\) 与 \(u_{2}\) 属于 \(\mathcal{L}(\mathrm{E})\)。设 \(\varrho\) 为映射 \(v \mapsto u_{1} \circ v \circ u_{2}\)，它是从 E 到 E 的 Hilbert--Schmidt 映射所成的空间 \(\mathcal{L}_{2}(\mathrm{E})\) 的一个自同态。它的谱等于 \(\mathrm{Sp}(u_{1})\mathrm{Sp}(u_{2}) \subset \mathbf{C}\)。（应用类似的方法。）

\begin{enumerate}
\def\labelenumi{\arabic{enumi})}
\setcounter{enumi}{9}
\item
  设 E 为复 Hilbert 空间。设 \(x_1\) 与 \(x_2\) 为 E 的自同态。证明 \(x_1 \otimes x_2\) 作为 E \(\widehat{\otimes}_2\) E 的自同态，其谱等于 \(\mathrm{Sp}(x_1) \mathrm{Sp}(x_2)\)。（把 \(\overline{\mathrm{E}} \widehat{\otimes}_2\) E 与从 E 到 E 的 Hilbert--Schmidt 算子所成的空间等同起来，参见 EVT 5，第 52 页，定理 1，b)，并应用前一个习题。）
\item
  设 \(u\) 为复 Hilbert 空间 \(E\) 的自同态，并设 \((j, |u|)\) 为它的极分解。
\end{enumerate}

\begin{enumerate}
\def\labelenumi{\alph{enumi})}
\item
  为使 \(j\) 与 \(|u|\) 交换，充要条件是 \(u\) 与 \(u^{*}u\) 交换。
\item
  证明此时 \(\|u^{*}(x)\|\leqslant\|u(x)\|\) 对一切 \(x\in E\) 成立，并且 \(\varrho(u)=\|u\|\)。
\item
  设 \(E = \ell^{2}(\mathbf{N})\)，并设 \(u\) 为由下式给出的 E 的自同态
\end{enumerate}

\[
u ((a _ {0}, a _ {1}, \dots)) = (0, a _ {0}, a _ {1}, \dots).
\]

它不是正规的。计算 u 的极分解 \((j, |u|)\)；\(j\) 与 \(|u|\) 交换。

\begin{enumerate}
\def\labelenumi{\arabic{enumi})}
\setcounter{enumi}{11}
\tightlist
\item
  设 \(u\) 为 Banach 空间 E 的一个自同态。用 \(\sigma_{a}(u)\) 表示满足下列条件的 \(\lambda\in\mathbf{C}\) 所成的集合：对每个 \(\alpha\in\mathbf{R}_{+}^{*}\)，存在 \(x\in\mathrm{E}\) 使 \(\|\lambda x-u(x)\|<\alpha\|x\|\)。
\end{enumerate}

\begin{enumerate}
\def\labelenumi{\alph{enumi})}
\item
  证明 \(\sigma_{a}(u)\) 是 \(\sigma(u)\) 的一个闭子集。
\item
  证明 \(\sigma_{a}(u)\) 包含 \(\sigma(u)\) 的边界。
\end{enumerate}

\begin{enumerate}
\def\labelenumi{\arabic{enumi})}
\setcounter{enumi}{12}
\tightlist
\item
  设 E 为复 Hilbert 空间。设 N 为从 \(\mathcal{L}(\mathrm{E})\) 到 \(\mathfrak{P}(\mathbf{C})\) 中的一个映射，使得
\end{enumerate}

\begin{enumerate}
\def\labelenumi{\alph{enumi})}
\item
  \(N(u)\) 对每个 \(u \in \mathcal{L}(\mathrm{E})\) 都是紧的；
\item
  \(\mathrm{N}(\alpha u + \beta 1_{\mathrm{E}}) = \alpha \mathrm{N}(u) + \beta\) 对一切 \(u \in \mathcal{L}(\mathrm{E})\) 与每个 \((\alpha, \beta) \in \mathbf{C}^{2}\) 成立；
\item
  \(\mathrm{N}(u)\subset\{z\in\mathbf{C}\mid\mathcal{R}(z)\geqslant0\}\) 当且仅当 \(u+u^{*}\) 是正的。
\end{enumerate}

  那么 \(\mathrm{N}(u)=\overline{\iota(u)}\) 对一切 \(u\in\mathcal{L}(\mathrm{E})\) 成立。（证明 \(u\mapsto\overline{\iota(u)}\) 具有所指出的性质；设 \(N_{1}\) 与 \(N_{2}\) 为满足这些性质的映射；假设 \(\mathrm{N}_{1}(u)\) 不包含于 \(\mathrm{N}_{2}(u)\)，求 C 中的 \(\alpha\) 与 \(\beta\)，使 \(\mathrm{N}_{2}(\alpha u+\beta)\) 包含于半平面 \(\mathcal{R}(z)\geqslant0\) 而 \(\mathrm{N}_{1}(\alpha u+\beta)\) 不然，再利用 c) 得出矛盾。）

\begin{enumerate}
\def\labelenumi{\arabic{enumi})}
\setcounter{enumi}{13}
\item
  给出一个复 Hilbert 空间及 \(u \in \mathcal{L}(\mathrm{E})\) 的例子，使 \(\iota(u)\) 不是闭的。
\item
  C 的非空紧子集所成的集合 X 赋以 TG，IX，第 91 页，习题 6 中定义的距离。设 E 为复 Hilbert 空间。
\end{enumerate}

\begin{enumerate}
\def\labelenumi{\alph{enumi})}
\item
  映射 \(u \mapsto \overline{\iota(u)}\) 从 \(\mathcal{L}(\mathrm{E})\) 到 X 是连续的。若 E 是无穷维的，则当 \(\mathcal{L}(\mathrm{E})\) 赋以单收敛拓扑时，它不连续。
\item
  映射 \(u \mapsto \operatorname{Sp}(u)\) 从 \(\mathcal{L}(\mathrm{E})\) 到 X 不连续。
\end{enumerate}

\begin{enumerate}
\def\labelenumi{\arabic{enumi})}
\setcounter{enumi}{15}
\tightlist
\item
  设 X 为局部紧空间，\(\mu\) 为 X 上总质量为 1 的正测度，\(f: X \to X\) 为满足 \(f(\mu) = \mu\) 的映射。
\end{enumerate}

\begin{enumerate}
\def\labelenumi{\alph{enumi})}
\item
  映射 \(\varphi \mapsto \varphi \circ f\) 诱导出酉自同态 \(u_{f}\)，它定义在 \(\mathrm{L}^2 (\mathrm{X},\mu)\) 上。
\item
  实数 1 是 \(u_f\) 的重数为 1 的特征值，其充要条件是：对 X 的每个满足 \(f(A) = A\) 的可测子集 A，有 \(\mu(A) \in \{0,1\}\)。这时称 \(f\) 是 \(\mu\)-遍历的。
\item
  设 \(\mathrm{E}\subset\mathrm{L}^{2}(\mathrm{X},\mu)\) 为 \(u_{f}\) 关于特征值 1 的特征子空间。序列 \((v_{\mathrm{N}})_{\mathrm{N}\in\mathbf{N}}\) 定义为
\end{enumerate}

\[
v _ {\mathrm{N}} = \frac {1}{\mathrm{N}} \sum_ {n = 0} ^ {\mathrm{N-1}} u _ {f} ^ {n}
\]

在赋以单收敛拓扑的空间 \(\mathcal{L}(\mathrm{L}^{2}(\mathrm{X},\mu))\) 中收敛到 E 上的正交投影（"von Neumann 遍历定理"）。（注意 E 的正交补是 \(u_{f}-1_{\mathrm{E}}\) 的像的闭包；然后验证 \(v_{\mathrm{N}}(\varphi)\to0\) 当 \(\varphi\in\mathrm{E}^{\circ}\) 时成立。）

\(d)\) 定义 \(\widetilde{u}_f\)，即 \(\mathrm{L}^1 (\mathrm{X},\mu)\) 的一个如 \(a)\) 中那样的自同态，并令

\[
\widetilde {v} _ {\mathrm{N}} = \frac {1}{\mathrm{N}} \sum_ {n = 0} ^ {\mathrm{N-1}} \widetilde {u} _ {f} ^ {n}.
\]

对每个 \(\varphi\in\mathrm{L}^{1}(\mathrm{X},\mu)\)，序列 \((\widetilde{v}_{\mathrm{N}}(\varphi))_{\mathrm{N}\in\mathbf{N}}\) 在 \(\mathrm{L}^{1}(\mathrm{X},\mu)\) 中收敛到元素 \(\widetilde{\varphi}\)，它属于 \(\widetilde{u}_{f}\) 关于特征值 1 的特征子空间。（先考虑 \(\varphi\) 有界的情形，并应用前一个问题。）

\begin{enumerate}
\def\labelenumi{\alph{enumi})}
\setcounter{enumi}{4}
\tightlist
\item
  设 I 为有限集，X = I \(^{Z}\) 赋以 I 上的测度 \(\mu_n\) 的乘积测度，其中 \(\mu_n(i) = 1/\text{Card(I)}\) 对一切 \(i \in \text{I}\) 成立。令 \(f((i_n)) = (i_{n+1})\) 对一切 \((i_n) \in \text{X}\) 成立。那么 \(f(\mu) = \mu\)，且 \(f\) 是 \(\mu\)-遍历的（“Bernoulli 移位”）。
\end{enumerate}

\(f)\) 设 \(\mathrm{X} = \mathbf{R} / \mathbf{Z}\)，\(\mu\) 为 \(\mathrm{X}\) 上的规范化 Haar 测度。设 \(a \in \mathrm{X}\)，令 \(f(x) = x + a\)。那么 \(f(\mu) = \mu\)，且 \(f\) 是 \(\mu\)-遍历的当且仅当 \(a\) 是某个无理数的类。

